\chapter{Dyslexia}

\section*{Definition}
Dyslexia is a specific learning disability of neurobiological origin characterized by difficulty with accurate and/or fluent word recognition, poor decoding, and impaired spelling abilities despite adequate intelligence, instruction, and socio-economic opportunity. It affects approximately 5--10\% of the population.

\section*{What Happens in Your Body}
Dyslexia arises from dysfunction in the left hemisphere reading network, including the temporoparietal (word analysis) and occipitotemporal (word form) regions. Functional MRI studies show reduced activation in these areas and compensatory overactivation in the right hemisphere and frontal lobes. The leading theory posits a phonological processing deficit: difficulty segmenting spoken words into constituent phonemes and mapping them to graphemes. Diffusion tensor imaging reveals altered white matter microstructure in the left arcuate fasciculus and superior longitudinal fasciculus. Genetic studies identify risk variants in genes related to neuronal migration and cilia function (DCDC2, KIAA0319, ROBO1).

\section*{Causes}
\begin{itemize}
  \item \textbf{Genetic:} Family history in 40--60\% of cases; multiple genes with small effect sizes
  \item \textbf{Neuroanatomic:} Ectopias and microgyria in left perisylvian cortex; reduced gray matter volume
  \item \textbf{Phonological deficit:} Core impairment in phoneme awareness and manipulation
  \item \textbf{Visual processing:} Magnocellular pathway dysfunction (less supported evidence)
  \item \textbf{Rapid auditory processing:} Difficulty perceiving rapid acoustic changes in speech
  \item \textbf{Comorbidities:} ADHD (25--40\%), dysgraphia, dyscalculia, language impairment
\end{itemize}

\section*{When to See a Doctor}
See your doctor if:
\begin{itemize}
  \item Child has difficulty learning letter names and sounds despite adequate instruction
  \item Reading accuracy or fluency is significantly below grade level (2+ grade levels)
  \item Child avoids reading aloud or demonstrates frustration with reading tasks
  \item Poor spelling persists despite phonetically predictable writing instruction
  \item Family history of reading difficulties or diagnosed dyslexia
\end{itemize}

\section*{Related Symptoms}
\begin{itemize}
  \item Slow, effortful, and inaccurate reading
  \item Difficulty segmenting and blending sounds (phonological awareness)
  \item Poor spelling with phonetic errors
  \item Trouble learning foreign languages
  \item Difficulty with rapid automatized naming (colors, objects, letters)
  \item Reduced working memory and verbal processing speed
  \item Concomitant anxiety, low self-esteem, and avoidance of reading
\end{itemize}
\textbf{Important:} Dyslexia is NOT a vision problem; vision therapy and colored lenses have no evidence of efficacy for the core phonological deficit.

\section*{Self-Care / Home Management}
\begin{itemize}
  \item Reading aloud with the child for 15--20 minutes daily
  \item Multisensory reading programs (e.g., Orton-Gillingham approach)
  \item Audiobooks and text-to-speech software as reading accommodations
  \item Consistent homework routines with frequent breaks
  \item Positive reinforcement focused on effort rather than reading performance
\end{itemize}

\section*{How Your Doctor Will Evaluate You}
\begin{itemize}
  \item Comprehensive psychoeducational evaluation: IQ testing (WISC-V), academic achievement testing (WIAT-IV, Woodcock-Johnson)
  \item Phonological processing assessment: CTOPP-2 (Comprehensive Test of Phonological Processing)
  \item Reading fluency and rapid naming tests (TOWRE-2)
  \item Spelling and writing assessment
  \item Vision and hearing screening to exclude sensory deficits
  \item Behavioral and attention screening to identify comorbid conditions
\end{itemize}

\section*{Treatment Options}
\begin{itemize}
  \item Structured literacy instruction with explicit phonics, phonemic awareness, and decoding
  \item Evidence-based intervention programs (Orton-Gillingham, Wilson Reading System, Lindamood-Bell)
  \item Recommendations for individualized education plan (IEP) or 504 plan accommodations
  \item Assistive technology: speech-to-text, text-to-speech, audiobooks
  \item Extended time on tests, oral examination alternatives, note-taking support
  \item Psychological support for anxiety and self-esteem issues
  \item Accommodations for standardized testing (extended time, screen reader)
\end{itemize}

\clearpage
